% ============================================================
% CHƯƠNG 4 — XÂY DỰNG HỆ THỐNG (Giai đoạn 3–4 của phiếu)
% Khung chương: đề mục + gợi ý trong chú thích + hình/bảng đã vẽ sẵn (figures_chap4.tex).
% \canviet{...} = chỗ bạn tự viết. Mã nguồn tương ứng nằm trong contracts/, backend/, frontend/.
% ============================================================

\chapter{XÂY DỰNG HỆ THỐNG}
\label{chap:xaydung}
% ============================================================
% HÌNH VÀ BẢNG CHƯƠNG 4 — gọi bằng macro trong Chap4.tex
%   \BangCongNghe   : công nghệ và phiên bản sử dụng      (tab:congnghe)
%   \HinhTrienKhai  : sơ đồ triển khai hệ thống            (fig:trienkhai)
%   \HinhCauTrucMa  : cấu trúc thư mục mã nguồn            (fig:cautrucma)
% Phiên bản lấy từ package.json / hardhat.config.js tại thời điểm 21/10/2026.
% ============================================================

\newcommand{\BangCongNghe}{%
\begin{table}[htbp]
\centering
\caption{Công nghệ và phiên bản sử dụng}
\label{tab:congnghe}
\small
\renewcommand{\arraystretch}{1.2}
\begin{tabularx}{\textwidth}{|p{3.2cm}|p{4.4cm}|X|}
\hline
\textbf{Thành phần} & \textbf{Công nghệ (phiên bản)} & \textbf{Vai trò}\\ \hline
Hợp đồng thông minh & Solidity 0.8.24, OpenZeppelin Contracts 5.1.0 & ERC-721, ERC-20, AccessControl, Pausable, ReentrancyGuard\\ \hline
Môi trường phát triển & Hardhat 2.29, hardhat-toolbox 5 (ethers 6, Mocha/Chai, solidity-coverage, gas reporter) & Biên dịch, kiểm thử, đo độ phủ và gas, triển khai\\ \hline
Mạng blockchain & Hardhat Network (cục bộ), Ethereum Sepolia & Kiểm thử và triển khai thử nghiệm\\ \hline
Backend & Node.js 22, Express 4.21, jsonwebtoken 9, multer & API REST, xác thực bằng chữ ký ví, tải hồ sơ\\ \hline
Cơ sở dữ liệu & PostgreSQL 16 (pg 8.13); pg-mem khi kiểm thử & Lưu dữ liệu off-chain theo mô hình CAD Trust\\ \hline
Đồng bộ sự kiện & ethers 6.13 (JsonRpcProvider, queryFilter) & Đọc sự kiện on-chain vào CSDL\\ \hline
Frontend & Next.js 14.2 (App Router), React 18.3, ethers 6.13 & Giao diện web, kết nối MetaMask\\ \hline
Ví & MetaMask & Ký đăng nhập và ký giao dịch\\ \hline
\end{tabularx}
\end{table}}

\newcommand{\HinhTrienKhai}{%
\begin{figure}[htbp]
\centering
\begin{adjustbox}{max width=\textwidth}
\begin{tikzpicture}[>=Stealth, font=\footnotesize,
    hop/.style={draw, rounded corners=3pt, align=center, fill=white, minimum height=1.1cm, text width=3.2cm},
    nhan/.style={font=\scriptsize, align=center}]
  \node[hop] (br)  at (0,0)    {\textbf{Trình duyệt}\\ giao diện Next.js\\ + MetaMask};
  \node[hop] (fe)  at (0,-3.2) {\textbf{Máy chủ frontend}\\ Next.js (Vercel)};
  \node[hop] (be)  at (5.2,0)  {\textbf{Máy chủ backend}\\ Express, port 4000\\ + bộ đồng bộ sự kiện};
  \node[hop] (db)  at (10.4,0) {\textbf{PostgreSQL}};
  \node[hop] (fs)  at (10.4,-3.2) {\textbf{Kho tệp}\\ uploads/ (hoặc IPFS)};
  \node[hop] (rpc) at (5.2,-3.2) {\textbf{Nút RPC}\\ (Infura/Alchemy)};
  \node[hop, fill=gray!10] (eth) at (5.2,-6.0) {\textbf{Ethereum Sepolia}\\ CarbonCreditNFT, CarbonPool};
  \draw[->] (fe) -- node[nhan, left] {tải trang\\ (HTTPS)} (br);
  \draw[<->] (br) -- node[nhan, above] {REST/JSON\\ + JWT} (be);
  \draw[<->] (be) -- node[nhan, above] {SQL} (db);
  \draw[->] (be) -- node[nhan, right] {lưu tệp} (fs);
  \draw[->] (be) -- node[nhan, right] {đọc sự kiện} (rpc);
  \draw[->] (br.south east) to[bend left=10] node[nhan, pos=0.35, right] {MetaMask ký\\ và gửi giao dịch} (rpc.north west);
  \draw[<->] (rpc) -- (eth);
\end{tikzpicture}
\end{adjustbox}
\caption{Sơ đồ triển khai hệ thống}
\label{fig:trienkhai}
\end{figure}}

\newcommand{\HinhCauTrucMa}{%
\begin{figure}[htbp]
\centering
\footnotesize
\renewcommand{\arraystretch}{1.05}
\begin{tabular}{@{}l l@{}}
\texttt{DATN-CarbonCreditOnBlockChain/} & \\
\texttt{\quad contracts/} & \textit{Hardhat}\\
\texttt{\quad\quad contracts/CarbonCreditNFT.sol} & lô tín chỉ ERC-721: mint, retire\\
\texttt{\quad\quad contracts/CarbonPool.sol} & pool ERC-20: deposit, redeem, retire\\
\texttt{\quad\quad test/CarbonCredit.test.js} & 26 test case TC-01 đến TC-26\\
\texttt{\quad\quad scripts/deploy.js} & triển khai, cấp quyền\\
\texttt{\quad\quad scripts/demo.js} & kịch bản demo, đo gas\\
\texttt{\quad\quad scripts/e2e-indexer.js} & kiểm thử on-chain $\rightarrow$ CSDL\\
\texttt{\quad\quad scripts/smoke-api.js} & kiểm thử toàn hệ thống đang chạy\\
\texttt{\quad backend/} & \textit{Node.js + Express}\\
\texttt{\quad\quad db/schema.sql} & lược đồ CSDL (CAD Trust + ứng dụng)\\
\texttt{\quad\quad src/routes/*.js} & các nhóm endpoint API\\
\texttt{\quad\quad src/indexer.js} & đồng bộ sự kiện on-chain\\
\texttt{\quad\quad test/api.test.js} & 16 test case API-01 đến API-16\\
\texttt{\quad frontend/} & \textit{Next.js}\\
\texttt{\quad\quad app/} & các trang giao diện\\
\texttt{\quad\quad components/WalletProvider.js} & đăng nhập bằng ví\\
\texttt{\quad\quad lib/chain.js, lib/api.js} & gọi hợp đồng và API\\
\end{tabular}
\caption{Cấu trúc thư mục mã nguồn}
\label{fig:cautrucma}
\end{figure}}


% Đoạn mở đầu: chương này cài đặt thiết kế ở Chương 3 theo ba phần (hợp đồng, backend, frontend)
% rồi triển khai lên testnet.
% [VIẾT SAU] \canviet{Đoạn mở đầu chương}

% ------------------------------------------------------------
\section{Môi trường và công nghệ}
\label{sec:moitruong}
% Giải thích \autoref{tab:congnghe}; lý do chọn Hardhat (test bằng JS, coverage, gas reporter),
% OpenZeppelin 5 (thư viện đã kiểm toán). Cấu trúc monorepo theo \autoref{fig:cautrucma}.
% [VIẾT SAU] \canviet{Môi trường và công nghệ}

\BangCongNghe

\HinhCauTrucMa

% ------------------------------------------------------------
\section{Cài đặt hợp đồng thông minh}
\label{sec:caidat-hopdong}

\subsection{Hợp đồng CarbonCreditNFT}
\label{subsec:caidat-nft}
% Ý chính (đối chiếu contracts/contracts/CarbonCreditNFT.sol):
%  - mỗi NFT = một lô (Unit): struct CreditBatch (projectId, vintage, serialFrom, quantity, retired,
%    documentHash, verifier, issuedAt);
%  - mintBatch: onlyRole(VERIFIER_ROLE); hai kiểm tra chống phát hành kép:
%    (1) seriesKey = keccak256(projectId, vintage), serialFrom >= nextAvailableSerial[key];
%    (2) documentUsed[documentHash] chưa dùng;
%  - retire: hủy một phần hoặc toàn bộ, luôn hủy từ đầu dãy sê-ri còn lại -> chứng nhận có dãy sê-ri
%    chính xác; hủy hết thì _burn NFT; dữ liệu lô vẫn tra cứu được qua getBatch;
%  - retireFor: dành cho pool (RETIREMENT_AGENT_ROLE) hủy thay người đã đốt token;
%  - Pausable: cơ quan quản lý tạm dừng toàn bộ mint/chuyển/hủy (ghi đè _update).
% Nên trích đoạn mã mintBatch và _retire (lstlisting ngắn) và giải thích từng kiểm tra.
% [VIẾT SAU] \canviet{Cài đặt CarbonCreditNFT}

\subsection{Hợp đồng CarbonPool}
\label{subsec:caidat-pool}
% Ý chính (CarbonPool.sol): decimals = 0 (1 PCC = 1 tCO2e); mọi đường nạp NFT đều qua
% onERC721Received (chỉ nhận NFT từ hợp đồng tín chỉ); redeem đốt token = remainingOf(tokenId);
% retire đốt token rồi gọi retireFor; bất biến totalSupply = tổng remainingOf các lô trong pool;
% ReentrancyGuard cho deposit/redeem/retire. So sánh với Toucan (BCT/TCO2): ở đây pool gắn trực
% tiếp với lô cụ thể khi hủy nên chứng nhận vẫn có dãy sê-ri.
% [VIẾT SAU] \canviet{Cài đặt CarbonPool}

\subsection{Phân quyền và triển khai}
\label{subsec:phanquyen}
% Bảng vai trò: DEFAULT_ADMIN_ROLE (cơ quan quản lý, ví deploy), VERIFIER_ROLE (VVB),
% RETIREMENT_AGENT_ROLE (hợp đồng pool). Các bước trong scripts/deploy.js. Địa chỉ hợp đồng
% trên Sepolia + link Etherscan (điền sau khi deploy thật).
% [VIẾT SAU] \canviet{Phân quyền và triển khai hợp đồng}

% ------------------------------------------------------------
\section{Cài đặt backend}
\label{sec:caidat-backend}

\subsection{Cơ sở dữ liệu}
\label{subsec:caidat-csdl}
% Từ ERD ở \autoref{fig:erd} sang db/schema.sql: nhóm bảng CAD Trust và nhóm bảng ứng dụng;
% cách tách khối Units khi hủy một phần (Held -> Held + Retired).
% [VIẾT SAU] \canviet{Cài đặt cơ sở dữ liệu}

\subsection{Xác thực bằng chữ ký ví}
\label{subsec:xacthuc}
% Luồng nonce -> ký thông điệp (personal_sign) -> ethers.verifyMessage -> JWT 8 giờ;
% nonce dùng một lần (chống phát lại); vai trò VERIFIER/ADMIN không tự đăng ký được.
% [VIẾT SAU] \canviet{Xác thực bằng chữ ký ví}

\subsection{API REST}
\label{subsec:caidat-api}
% Giải thích \autoref{tab:api} theo nhóm: xác thực, dự án, hồ sơ, yêu cầu cấp, tra cứu, dashboard,
% liên thông CAD Trust. Nhấn mạnh: backend không giữ khóa riêng, chỉ trả tham số mint cho VVB ký.
% [VIẾT SAU] \canviet{API REST}

\BangAPI

\subsection{Bộ đồng bộ sự kiện}
\label{subsec:indexer}
% src/indexer.js: quét theo đoạn khối (queryFilter), sắp xếp theo (block, logIndex), mỗi sự kiện
% ghi trong một giao dịch CSDL, khóa (tx_hash, log_index) để không ghi trùng; on-chain là gốc,
% có thể dựng lại CSDL từ khối deploy. Hạn chế: chưa xử lý reorg (Sepolia cần đợi vài khối xác nhận).
% [VIẾT SAU] \canviet{Bộ đồng bộ sự kiện}

% ------------------------------------------------------------
\section{Cài đặt giao diện}
\label{sec:caidat-frontend}
% Mỗi màn hình 1 ảnh chụp + 2–3 câu: trang chủ/danh sách dự án, chi tiết dự án, tra cứu sê-ri,
% chứng nhận hủy, dự án của tôi (chủ dự án), thẩm định (VVB), ví tín chỉ, dashboard, quản trị.
% CHỤP ẢNH khi chạy thật (xem README mục "Chạy thử cục bộ").
% [VIẾT SAU] \canviet{Các màn hình chính kèm ảnh chụp}

% ------------------------------------------------------------
\section{Triển khai hệ thống}
\label{sec:trienkhai}
% Giải thích \autoref{fig:trienkhai}: Sepolia + RPC, backend + PostgreSQL, frontend.
% [VIẾT SAU] \canviet{Triển khai hệ thống}

\HinhTrienKhai

\section{Kết luận chương}
% [VIẾT SAU] \canviet{Kết luận chương 4}
